\chapter{Future Work}
\label{chapter:futurework}

Overall, this thesis highlights an important segment of the broader topic of LLMs in software engineering, while also indicating directions for future research.
In Chapter~\ref{chapter_2}, optional comparison criteria could be incorporated, and an automated system might be developed to perform functional and non-functional tests.
While human evaluation cannot be fully automated, a substantial portion of the process could be streamlined.

In Chapter~\ref{chapter_4}, the patterns associated with the temperature hyper parameter can be further explored at multiple levels.
This includes not only understanding its effects on model behavior but also determining the optimal granularity for its application.

Although Chapter~\ref{chapter_3} may appear complete, given that GPT-4 is now considered an older model, there remains room for further work.
Newer models could be evaluated, and improved inference techniques could be developed to achieve higher success rates.
While the one-third success rate observed is promising, real-world developers require even more reliable outcomes.

\vspace{5em}
To conclude this thesis, Large Language Models have become an integral part of software engineering, but they still require extensive research across multiple areas and levels to fully meet the needs of the developer community and the developer community needs to understand and adapt to these new tools.

